% Root and B full literal audits passed; D's finite counterexample gate passed.
\subsection{Natural graph comparisons for every opposite corner of order at most four}

\begin{lemma}[A natural graph for an identity-adjoined right-zero band]
\label{C:right-zero-band-natural-graph}
Let $N=\{1\}\sqcup K$, where $K$ is a nonempty finite right-zero
band, and let $S$ be a finite left $N$-set.  Put $U=KS$.
There is a natural finite graph $G_N(S)$ and an actual weak
equivalence
\[
 B(N\ltimes S)\longrightarrow G_N(S).
\]
Its vertex set is $S$.  For each $s\notin U$ and $k\in K$,
it has a labeled edge from $s$ to $ks$.  All maps of left
$N$-sets induce graph maps taking an edge at unit speed to an
edge with the same label or collapsing it to a vertex.
In particular the comparison is equivariant for every commuting
right monoid action.
\end{lemma}
\begin{proof}
Every $k\in K$ fixes $U$ pointwise.  Regard the displayed graph
as a category with the indicated edges as its only nonidentity
arrows; there are no composable nonidentity pairs.  Define the
functor from the action category by keeping its objects and
its arrows from $S\setminus U$ to $U$, and sending every
endomorphism at a core object to its identity.  The right-zero
multiplication makes this a functor.

For a core vertex $u$, the comma category $(u\downarrow q)$
has one object and endomorphism monoid $N$, whose classifying
space is contractible.  For a noncore vertex $s$, the comma
category $(s\downarrow q)$ has the initial object $(s,1_s)$.
Indeed the other objects are indexed by the individual labeled
edges from $s$, and the corresponding labeled action arrow is
the unique morphism from $(s,1_s)$ to each of them.  Quillen's
Theorem A proves the weak equivalence.  A map of acts preserves
the core, and sends a noncore vertex either to another noncore
vertex or to the core.  The stated edge rule therefore makes
both the graph construction and the functor strictly natural.
\end{proof}

\begin{lemma}[The higher homology of a two-band four-element action category]
\label{C:four-corner-action-tail}
Let $N=\{1,p,q,x\}$ have minimal ideal the right-zero band
$K=\{p,q\}$.  Let $S$ be a finite left $N$-set, $U=KS$,
$R=\mathbb QN$, and $A=R/\mathbb QK$.  Write
$W=\mathbb Q(q-p)$ as a right $A$-module and
$T=\mathbb QS/\mathbb QU$ as a left $A$-module.  Then
\[
 H_n(B(N\ltimes S);\mathbb Q)
   \cong \operatorname{Tor}_{n-1}^{A}(W,T),
 \qquad n\geq2.
\]
The comparison is natural for every commuting right action
on $S$.  In four of the five such four-element monoids, $A$
is semisimple, so these homology groups vanish.  The sole
exception, up to relabeling $p,q$, has
\[
 px=qx=p,\qquad xp=p,\qquad xq=q,\qquad x^2=p.
\]
For this exception $A=\mathbb Q[\epsilon]/(\epsilon^2)$
and $W$ is its residue module.
\end{lemma}
\begin{proof}
Put $e=p$.  Then $Re=\mathbb Qp$ and
$eR=\mathbb QK=\mathbb Qp\oplus\mathbb Qq$.
The multiplication isomorphism
$\mathbb QK=Re\otimes_{\mathbb Q}eR$ gives
\[
 \mathbb QK\otimes_R\mathbb QS
   =Re\otimes_{\mathbb Q}e\mathbb QS
   \cong\mathbb QU.
\]
Its map to $\mathbb QS$ is the literal inclusion of the
core-state span.  Since
$0\to\mathbb QK\to R\to A\to0$ is a projective resolution
of the right $R$-module $A$, this proves the natural derived
identification
\[
 A\otimes_R^{\mathbb L}\mathbb QS\simeq T.
\]
The right augmentation sequence
$0\to W\to eR\to\mathbb Q\to0$ has projective middle term.
Dimension shifting and derived associativity consequently give
\[
 H_n(B(N\ltimes S);\mathbb Q)
 =\operatorname{Tor}_{n}^{R}(\mathbb Q,\mathbb QS)
 =\operatorname{Tor}_{n-1}^{R}(W,\mathbb QS)
 =\operatorname{Tor}_{n-1}^{A}(W,T)
\]
in exactly the stated range.

For completeness the five possibilities are as follows.
If the right translation of $x$ on $K$ is the identity, then
$x^2=1$ or $x$.  If it swaps $p,q$, then $x^2=1$.
If it is constant, relabeling makes that constant $p$,
and $x^2=x$ or $p$.  These possibilities follow by applying
the right translation to $x^2$, and excluding a larger
minimal ideal.  The first four contracted algebras have
basis $1,\bar x$ with $\bar x^2=1$ or $\bar x$, hence
are semisimple over $\mathbb Q$.  The final possibility
has $\bar x^2=0$.  There the column-difference character
annihilates $\bar x$, as asserted.
\end{proof}

\begin{proposition}[The full natural small-corner comparison]
\label{C:small-opposite-natural-graph-comparison}
Let $|N|\leq4$, let $M$ be any finite monoid, and let $S$
be a nonempty finite $(N,M)$-biset.  There is a finite graph
$\Gamma$, with a right $M$-action by unit-speed edge maps
(allowing reversal) and edge collapses, and a natural
comparison from the collage to
\[
 Y=\operatorname{hocofib}(\Gamma_{hM}\longrightarrow BM),
\]
through actual maps and equivariant weak-equivalence zigzags,
which induces
\[
 H^2(Y;\mathbb Q)\xrightarrow{\;\cong\;}H^2(BC;\mathbb Q).
\]
Thus products on the collage are pullbacks of products on
the graph cofiber.  If $N$ is not the exceptional nilpotent
four-element monoid of Lemma~\ref{C:four-corner-action-tail},
this comparison is a rational homology equivalence in every
degree.  For that exception only the displayed low-degree
assertion is made.
\end{proposition}
\begin{proof}
If the trivial right $\mathbb QN$-module is projective,
the actual action-category projection
$N\ltimes S\to N\backslash S$ is a rational homology
equivalence: its chain homology is computed by the exact
coinvariant functor.  Also $BN$ is rationally acyclic.
Take $\Gamma$ to be the discrete right act $N\backslash S$.
The natural collage quotient, or its double-bar homotopy
pushout, gives the asserted full rational comparison.

The elementary small-monoid inventory leaves, in the
nonprojective case, $1+$ a right-zero two-band, the five
two-band four-element monoids, and $1+$ a right-zero
three-band.  The identity-adjoined bands have the natural
actual graph model of
Lemma~\ref{C:right-zero-band-natural-graph}.  Their classifying
spaces are contractible, so their collages are the indicated
graph cofibers up to homotopy.

For the five remaining monoids use the normalized
groupoid presentation of
Lemma~\ref{C:four-corner-normalized-action-groupoid}, for
arbitrary $S$, with no state bound.  Its objects are the
core $U$, and the generators are $g_s:Ps\to Qs$.
When the right translation of $x$ is the identity, quotient
the parallel edges by $g_{Es}=g_s$, collapsing an edge when
its image is a core vertex.  When the translation is constant,
collapse the loops $g_{Es}$; this is an invariant set.
Both constructions give finite oriented graphs with strict
right actions by edge maps or collapses.

In the swap case $E^2=1$.  Each two-element outside orbit
gives one unoriented edge.  An $E$-fixed outside point gives
a loop with relation $g_s^2=1$; kill these order-two loop
generators.  The result is the fundamental groupoid of a
finite graph.  Commuting right maps preserve the fixed
points and the pairs, so they give strict graph maps,
possibly reversing an edge.  Killing the order-two
generators preserves first rational homology.

In all three cases the actual normalized functor from
$N\ltimes S$ to the resulting graph groupoid is natural
and induces a component bijection and an isomorphism on
first rational homology.  The natural one-truncation
equivalence of a graph with the classifying space of its
fundamental groupoid is an equivalence of $M$-diagrams.
One may realize this equivalence using singular simplicial
sets and their natural one-truncations; it does not require
choosing equivariant spanning trees or edge representatives.

The target graph has no second homology.  Consequently the
cone of the augmented chain comparison has zero homology
in degrees zero, one and two.  Tensoring with the free left
monoid bar preserves that connectivity.  Suspending the
augmented bar comparison gives an isomorphism on second
homology of the actual cofibers, and hence the displayed
second-cohomology isomorphism.  The cofiber maps are actual
maps of spaces, so cup products pull back.

In the four semisimple cases,
Lemma~\ref{C:four-corner-action-tail} also shows that the
source action category has no rational homology above one.
Its graph comparison is therefore a rational homology
equivalence in all degrees.  The free bar and the homotopy
pushout preserve that equivalence.  No higher-degree
equivalence is asserted in the remaining nilpotent case.
\end{proof}

\begin{remark}[A necessary restriction on the graph comparison]
\label{C:small-opposite-graph-higher-homology-warning}
The exceptional qualification is essential.  In the
nilpotent four-element monoid take $S=\{u,t\}$, with every
nonidentity actor sending both states to $u$.  Then
$T=\mathbb Q\{t\}$ and $W$ are both the residue module of
$A=\mathbb Q[\epsilon]/(\epsilon^2)$.  The periodic free
resolution with every differential multiplication by
$\epsilon$ gives
\[
 H_n(B(N\ltimes S);\mathbb Q)=\mathbb Q
 \quad\hbox{for every }n\geq2.
\]
The normalized graph is one circle; the same augmentation
sequence, or its fundamental-groupoid presentation, gives
$H_1=\mathbb Q$ as well.  Thus this finite action category
has nonzero rational homology in every positive degree.
Its natural graph comparison still induces the stated
second-cohomology isomorphism on collages, but cannot
transfer bounded higher homology from a collage to its
graph replacement.  No spectral-sequence degeneration
or arbitrary finite-type inference is being used.

The standalone program
\path{code/d_verify_nilpotent_graph_counterexample.py}
checks all $64$ multiplication triples and $32$ action
equations, the literal projective right ideals with bases
$(p,q)$ and $(1-p,x-p)$, and every kernel and image in the
periodic resolution.  It also rebuilds the eight-arrow
action category and verifies the actual functor to the
one-circle groupoid, together with a first-homology
cycle--cocycle pairing equal to one.  Its complete record,
\path{code/outputs/d_nilpotent_graph_counterexample_verified.json},
certifies all the stated finite data and the all-degree
periodic conclusion without producer imports.
\end{remark}

\begin{corollary}[Unconditional small-corner graph cup consequences]
\label{C:small-opposite-P5-all-states}
For every $|N|\leq4$ and every finite $(N,P_5)$-biset,
all rational degree-two products on the collage vanish,
with no state bound or bounded-homology assumption.
The same conclusion holds for a source $M$ whose minimal
ideal is a two-column rectangular band and whose row
augmentation is projective over $\mathbb QM/\mathbb QK$.
In particular the seventeen-arrow profile $(5,4,8)$ is
excluded in full.
\end{corollary}
\begin{proof}
Apply the second-cohomology isomorphism in the proposition.
The unconditional projective-row clause of
Theorem~\ref{C:two-column-graph-localization} kills all
degree-two products on the graph cofiber, including graphs
with reversed edges.  For $P_5$ its contracted algebra is
$\mathbb Q$, so its row augmentation is projective.
The empty-biset case is a disjoint union of the two vertex
spaces, whose fourth rational cohomology is zero.
\end{proof}

\begin{corollary}[The bounded small-source consequence away from the exception]
\label{C:small-opposite-six-source-regular}
Suppose $|M|\leq6$, $|N|\leq4$, and $N$ is not the
exceptional nilpotent four-element corner.  If the collage
of a finite $(N,M)$-biset has bounded rational homology,
all its rational degree-two products vanish, at every
state size.
\end{corollary}
\begin{proof}
The natural graph comparison is now a full rational
homology equivalence, so its graph cofiber has bounded
rational homology.  If $H^2(BM;\mathbb Q)=0$, the vertex
restrictions vanish and the suspension quotient kills
products.  Otherwise the minimal ideal of $M$ is a
$2\times2$ band and the contracted algebra has dimension
at most two.  The projective or local alternatives in
Theorem~\ref{C:two-column-graph-localization} apply exactly
as in Corollary~\ref{C:six-source-right-zero-all-states}.
The exceptional $N$ is deliberately excluded here.
\end{proof}

\begin{corollary}[The corresponding star comparisons]
\label{C:small-opposite-graph-stars}
Let a finite star category have central endomorphism
monoid $M$, sink endomorphism monoids $N_i$ of order at
most four, and nonempty cross bisets $S_i$.  If $M=P_5$,
or if $M$ has a two-column minimal band with projective
contracted row augmentation, all rational degree-two
products on its classifying space vanish.  If $|M|\leq6$,
the same conclusion holds under bounded rational
homology provided none of the $N_i$ is the exceptional
nilpotent corner.  All assertions have opposite versions.
\end{corollary}
\begin{proof}
Apply the natural comparison separately to each edge
bar in the star's homotopy-pushout description.  The
comparisons fix the central $BM$, and each $BN_i$ is
rationally acyclic.  They therefore give the same
second-cohomology isomorphism, and in the nonexceptional
case a full rational homology equivalence, with the
star of graph cofibers.

Identify its finitely many cone vertices.  For a connected
space, identifying $r$ distinct points adds a wedge of
$r-1$ circles up to homotopy.  The resulting space is
the single cofiber associated to the disjoint union
of the coefficient graphs.  This operation preserves
the cohomology ring in degrees at least two and preserves
bounded rational homology.  The graph-localization
theorem applies to that disjoint union and proves
the asserted conclusions.  The only boundedness transfer
used is the full rational comparison in the explicitly
nonexceptional case.
\end{proof}
